% --- Archon physics formalization source begin ---
% archon:physics
% archon:covers PhyXMiniProblems/problem_phyx_mini_0428.lean
% archon:source-report reports/phyx_mini/problem_phyx_mini_0428.source.json

\chapter{Physics problem phyx\_mini\_0428}
\label{ch:PhyXMiniProblems_problem_phyx_mini_0428}

\paragraph{Problem source.}
\#\# Physical scenario

A heat engine using a monatomic gas follows the cycle shown in figure.

\#\# Question

What is the thermal efficiency of this heat engine?

\#\# Answer choices

- A: A: 0.40

- B: B: 0.64

- C: C: 0.316

- D: D: 0.32

\#\# Figure caption (auxiliary; use the image as primary evidence)

The image is a pressure-volume (p-V) diagram representing a thermodynamic process. Here's a detailed breakdown:

- **Axes**: 
  - The vertical axis represents pressure (p) in kilopascals (kPa).
  - The horizontal axis represents volume (V) in cubic centimeters (cm³).

- **Process**:
  - The diagram shows a three-step cycle involving points 1, 2, and 3.
  - **1 to 2**: Horizontal arrow, indicating an isobaric (constant pressure) process at \textbackslash{}( p\_\{\textbackslash{}text\{max\}\} \textbackslash{}) with the temperature noted as 600 K.
  - **2 to 3**: Vertical arrow, indicating an isochoric (constant volume) process.
  - **3 to 1**: Curved arrow labeled "Adiabatic," indicating an adiabatic process (no heat exchange) returning to the initial point.

- **Calibration**:
  - Pressure ranges from 0 to just above 100 kPa.
  - Volume ranges from 0 to over 600 cm³.

- **Text and Labels**:
  - \textbackslash{}( p\_\{\textbackslash{}text\{max\}\} \textbackslash{}) near point 1 indicates maximum pressure.
  - A label denotes the process from 1 to 2 as happening at 600 K.
  - The curved process from 3 to 1 is labeled "Adiabatic." 

The diagram visualizes the relationships and transitions between pressure and volume in a thermodynamic cycle.

\#\# Dataset metadata

Laws of Thermodynamics; Physical Model Grounding Reasoning, Multi-Formula Reasoning

\paragraph{Recorded answer/context.}
D: D: 0.32

\paragraph{Figure/image path.}
/root/proposal\_for\_physic/science-mango/phyx\_mini\_run/phyx\_data/test\_image/428.png

\paragraph{Formalization target.}
create a compiling Lean file with sorry bodies at `PhyXMiniProblems/problem\_phyx\_mini\_0428.lean`. The Lean declarations must preserve the physical quantities, dimensions or dimensional roles, figure labels, governing-law hypotheses, and final relation expressed by this problem.
Use Mathlib/Physlib names found through LeanExplore where available. If a physics API is missing, introduce faithful local abstractions rather than scalar placeholder aliases.

\begin{theorem}[Physics formalization target]
\label{thm:physics:phyx_mini_0428:target}
\lean{PhyXMiniProblems.ProblemPhyXMini0428.problem_phyx_mini_0428}
\uses{def:physics:phyx-mini-0428:phyxminiproblems-problemphyxmini0428-monatomicidealgasheatenginecycle, def:physics:phyx-mini-0428:phyxminiproblems-problemphyxmini0428-matchesproblemstatement, def:physics:phyx-mini-0428:phyxminiproblems-problemphyxmini0428-matchessuppliedpressurevolumediagram, def:physics:phyx-mini-0428:phyxminiproblems-problemphyxmini0428-hasphysicalthermodynamicparameters, def:physics:phyx-mini-0428:phyxminiproblems-problemphyxmini0428-obeysmonatomicidealgascyclelaws, def:physics:phyx-mini-0428:phyxminiproblems-problemphyxmini0428-thermalefficiency, def:physics:phyx-mini-0428:phyxminiproblems-problemphyxmini0428-displayedefficiency, def:physics:phyx-mini-0428:phyxminiproblems-problemphyxmini0428-recordedanswerchoice, def:physics:phyx-mini-0428:phyxminiproblems-problemphyxmini0428-roundstotwodecimalplaces, def:physics:phyx-mini-0428:phyxminiproblems-problemphyxmini0428-exactefficiencyfromfigure}
The assigned autoformalize agent should translate this physics problem into Lean declarations in the covered file, with theorem and lemma proof bodies written as `by sorry`.
\end{theorem}
\begin{proof}
This is an autoformalization task, not a proof task. Produce statements that can later be proved by the physics prover without weakening the physical model.
\end{proof}

% --- Archon declaration topology begin ---
\section{Declaration topology}
\begin{definition}[Gas Volume]
  \label{def:physics:phyx-mini-0428:phyxminiproblems-problemphyxmini0428-gasvolume}
  \lean{PhyXMiniProblems.ProblemPhyXMini0428.GasVolume}
  A physical gas volume carrying the dimension length cubed
\end{definition}
\begin{proof}
  This is the declared model definition of Gas Volume.
\end{proof}
\begin{definition}[pressure In Pascals]
  \label{def:physics:phyx-mini-0428:phyxminiproblems-problemphyxmini0428-pressureinpascals}
  \lean{PhyXMiniProblems.ProblemPhyXMini0428.pressureInPascals}
  Coherent-SI pressure readout, in pascals
\end{definition}
\begin{proof}
  This is the declared model definition of pressure In Pascals.
\end{proof}
\begin{definition}[pressure In Kilopascals]
  \label{def:physics:phyx-mini-0428:phyxminiproblems-problemphyxmini0428-pressureinkilopascals}
  \lean{PhyXMiniProblems.ProblemPhyXMini0428.pressureInKilopascals}
  \uses{def:physics:phyx-mini-0428:phyxminiproblems-problemphyxmini0428-pressureinpascals}
  Pressure readout in the kilopascals printed on the vertical axis
\end{definition}
\begin{proof}
  This is the declared model definition of pressure In Kilopascals.
\end{proof}
\begin{definition}[volume In Cubic Metres]
  \label{def:physics:phyx-mini-0428:phyxminiproblems-problemphyxmini0428-volumeincubicmetres}
  \lean{PhyXMiniProblems.ProblemPhyXMini0428.volumeInCubicMetres}
  \uses{def:physics:phyx-mini-0428:phyxminiproblems-problemphyxmini0428-gasvolume}
  Coherent-SI volume readout, in cubic metres
\end{definition}
\begin{proof}
  This is the declared model definition of volume In Cubic Metres.
\end{proof}
\begin{definition}[volume In Cubic Centimetres]
  \label{def:physics:phyx-mini-0428:phyxminiproblems-problemphyxmini0428-volumeincubiccentimetres}
  \lean{PhyXMiniProblems.ProblemPhyXMini0428.volumeInCubicCentimetres}
  \uses{def:physics:phyx-mini-0428:phyxminiproblems-problemphyxmini0428-gasvolume, def:physics:phyx-mini-0428:phyxminiproblems-problemphyxmini0428-volumeincubicmetres}
  Volume readout in the cubic centimetres printed on the horizontal axis
\end{definition}
\begin{proof}
  This is the declared model definition of volume In Cubic Centimetres.
\end{proof}
\begin{definition}[energy In Joules]
  \label{def:physics:phyx-mini-0428:phyxminiproblems-problemphyxmini0428-energyinjoules}
  \lean{PhyXMiniProblems.ProblemPhyXMini0428.energyInJoules}
  Coherent-SI signed-energy readout, in joules
\end{definition}
\begin{proof}
  This is the declared model definition of energy In Joules.
\end{proof}
\begin{definition}[temperature In Kelvins]
  \label{def:physics:phyx-mini-0428:phyxminiproblems-problemphyxmini0428-temperatureinkelvins}
  \lean{PhyXMiniProblems.ProblemPhyXMini0428.temperatureInKelvins}
  Absolute-temperature readout in kelvins. The Temperature values in this setup are calibrated in the kelvin unit, as fixed by the 600 K label in the supplied figure
\end{definition}
\begin{proof}
  This is the declared model definition of temperature In Kelvins.
\end{proof}
\begin{definition}[Gas Model]
  \label{def:physics:phyx-mini-0428:phyxminiproblems-problemphyxmini0428-gasmodel}
  \lean{PhyXMiniProblems.ProblemPhyXMini0428.GasModel}
  The working-substance model stated in the problem
\end{definition}
\begin{proof}
  This is the declared model definition of Gas Model.
\end{proof}
\begin{definition}[Thermodynamic Device Role]
  \label{def:physics:phyx-mini-0428:phyxminiproblems-problemphyxmini0428-thermodynamicdevicerole}
  \lean{PhyXMiniProblems.ProblemPhyXMini0428.ThermodynamicDeviceRole}
  The thermodynamic role of the cyclic device
\end{definition}
\begin{proof}
  This is the declared model definition of Thermodynamic Device Role.
\end{proof}
\begin{definition}[State Label]
  \label{def:physics:phyx-mini-0428:phyxminiproblems-problemphyxmini0428-statelabel}
  \lean{PhyXMiniProblems.ProblemPhyXMini0428.StateLabel}
  Labels of the three black equilibrium states in the diagram
\end{definition}
\begin{proof}
  This is the declared model definition of State Label.
\end{proof}
\begin{definition}[Cycle Leg]
  \label{def:physics:phyx-mini-0428:phyxminiproblems-problemphyxmini0428-cycleleg}
  \lean{PhyXMiniProblems.ProblemPhyXMini0428.CycleLeg}
  Directed legs indicated by the arrowheads in the diagram
\end{definition}
\begin{proof}
  This is the declared model definition of Cycle Leg.
\end{proof}
\begin{definition}[Process Kind]
  \label{def:physics:phyx-mini-0428:phyxminiproblems-problemphyxmini0428-processkind}
  \lean{PhyXMiniProblems.ProblemPhyXMini0428.ProcessKind}
  Thermodynamic constraint written or geometrically displayed on a leg
\end{definition}
\begin{proof}
  This is the declared model definition of Process Kind.
\end{proof}
\begin{definition}[Path Geometry]
  \label{def:physics:phyx-mini-0428:phyxminiproblems-problemphyxmini0428-pathgeometry}
  \lean{PhyXMiniProblems.ProblemPhyXMini0428.PathGeometry}
  Geometric appearance of a leg in the primary raster
\end{definition}
\begin{proof}
  This is the declared model definition of Path Geometry.
\end{proof}
\begin{definition}[Axis Quantity]
  \label{def:physics:phyx-mini-0428:phyxminiproblems-problemphyxmini0428-axisquantity}
  \lean{PhyXMiniProblems.ProblemPhyXMini0428.AxisQuantity}
  Quantity assigned to one of the two diagram axes
\end{definition}
\begin{proof}
  This is the declared model definition of Axis Quantity.
\end{proof}
\begin{definition}[Axis Unit]
  \label{def:physics:phyx-mini-0428:phyxminiproblems-problemphyxmini0428-axisunit}
  \lean{PhyXMiniProblems.ProblemPhyXMini0428.AxisUnit}
  Unit printed next to a diagram axis
\end{definition}
\begin{proof}
  This is the declared model definition of Axis Unit.
\end{proof}
\begin{definition}[Diagram Label]
  \label{def:physics:phyx-mini-0428:phyxminiproblems-problemphyxmini0428-diagramlabel}
  \lean{PhyXMiniProblems.ProblemPhyXMini0428.DiagramLabel}
  Text or numerical annotations visibly present in the source image
\end{definition}
\begin{proof}
  This is the declared model definition of Diagram Label.
\end{proof}
\begin{definition}[expected Leg Start]
  \label{def:physics:phyx-mini-0428:phyxminiproblems-problemphyxmini0428-expectedlegstart}
  \lean{PhyXMiniProblems.ProblemPhyXMini0428.expectedLegStart}
  \uses{def:physics:phyx-mini-0428:phyxminiproblems-problemphyxmini0428-statelabel, def:physics:phyx-mini-0428:phyxminiproblems-problemphyxmini0428-cycleleg}
  Initial endpoint prescribed by the displayed arrow direction
\end{definition}
\begin{proof}
  This is the declared model definition of expected Leg Start.
\end{proof}
\begin{definition}[expected Leg Finish]
  \label{def:physics:phyx-mini-0428:phyxminiproblems-problemphyxmini0428-expectedlegfinish}
  \lean{PhyXMiniProblems.ProblemPhyXMini0428.expectedLegFinish}
  \uses{def:physics:phyx-mini-0428:phyxminiproblems-problemphyxmini0428-statelabel, def:physics:phyx-mini-0428:phyxminiproblems-problemphyxmini0428-cycleleg}
  Final endpoint prescribed by the displayed arrow direction
\end{definition}
\begin{proof}
  This is the declared model definition of expected Leg Finish.
\end{proof}
\begin{definition}[expected Process Kind]
  \label{def:physics:phyx-mini-0428:phyxminiproblems-problemphyxmini0428-expectedprocesskind}
  \lean{PhyXMiniProblems.ProblemPhyXMini0428.expectedProcessKind}
  \uses{def:physics:phyx-mini-0428:phyxminiproblems-problemphyxmini0428-cycleleg, def:physics:phyx-mini-0428:phyxminiproblems-problemphyxmini0428-processkind}
  Process classification read from the horizontal, vertical, and labelled curved paths
\end{definition}
\begin{proof}
  This is the declared model definition of expected Process Kind.
\end{proof}
\begin{definition}[expected Path Geometry]
  \label{def:physics:phyx-mini-0428:phyxminiproblems-problemphyxmini0428-expectedpathgeometry}
  \lean{PhyXMiniProblems.ProblemPhyXMini0428.expectedPathGeometry}
  \uses{def:physics:phyx-mini-0428:phyxminiproblems-problemphyxmini0428-cycleleg, def:physics:phyx-mini-0428:phyxminiproblems-problemphyxmini0428-pathgeometry}
  Geometry of each path in the pressure--volume plane
\end{definition}
\begin{proof}
  This is the declared model definition of expected Path Geometry.
\end{proof}
\begin{definition}[Thermodynamic State]
  \label{def:physics:phyx-mini-0428:phyxminiproblems-problemphyxmini0428-thermodynamicstate}
  \lean{PhyXMiniProblems.ProblemPhyXMini0428.ThermodynamicState}
  \uses{def:physics:phyx-mini-0428:phyxminiproblems-problemphyxmini0428-gasvolume}
  Pressure, volume, and absolute temperature at an equilibrium state
\end{definition}
\begin{proof}
  This is the declared model definition of Thermodynamic State.
\end{proof}
\begin{definition}[Pressure Volume Diagram]
  \label{def:physics:phyx-mini-0428:phyxminiproblems-problemphyxmini0428-pressurevolumediagram}
  \lean{PhyXMiniProblems.ProblemPhyXMini0428.PressureVolumeDiagram}
  \uses{def:physics:phyx-mini-0428:phyxminiproblems-problemphyxmini0428-statelabel, def:physics:phyx-mini-0428:phyxminiproblems-problemphyxmini0428-cycleleg, def:physics:phyx-mini-0428:phyxminiproblems-problemphyxmini0428-processkind, def:physics:phyx-mini-0428:phyxminiproblems-problemphyxmini0428-pathgeometry, def:physics:phyx-mini-0428:phyxminiproblems-problemphyxmini0428-axisquantity, def:physics:phyx-mini-0428:phyxminiproblems-problemphyxmini0428-axisunit, def:physics:phyx-mini-0428:phyxminiproblems-problemphyxmini0428-diagramlabel}
  Diagram-level data independent of the thermodynamic laws
\end{definition}
\begin{proof}
  This is the declared model definition of Pressure Volume Diagram.
\end{proof}
\begin{definition}[Monatomic Ideal Gas Heat Engine Cycle]
  \label{def:physics:phyx-mini-0428:phyxminiproblems-problemphyxmini0428-monatomicidealgasheatenginecycle}
  \lean{PhyXMiniProblems.ProblemPhyXMini0428.MonatomicIdealGasHeatEngineCycle}
  \uses{def:physics:phyx-mini-0428:phyxminiproblems-problemphyxmini0428-gasmodel, def:physics:phyx-mini-0428:phyxminiproblems-problemphyxmini0428-thermodynamicdevicerole, def:physics:phyx-mini-0428:phyxminiproblems-problemphyxmini0428-statelabel, def:physics:phyx-mini-0428:phyxminiproblems-problemphyxmini0428-cycleleg, def:physics:phyx-mini-0428:phyxminiproblems-problemphyxmini0428-thermodynamicstate, def:physics:phyx-mini-0428:phyxminiproblems-problemphyxmini0428-pressurevolumediagram}
  The physical gas sample, its state variables, and the signed energy transfers on the three legs. Neither thermal efficiency nor any answer choice is a field of this structure
\end{definition}
\begin{proof}
  This is the declared model definition of Monatomic Ideal Gas Heat Engine Cycle.
\end{proof}
\begin{definition}[Matches Problem Statement]
  \label{def:physics:phyx-mini-0428:phyxminiproblems-problemphyxmini0428-matchesproblemstatement}
  \lean{PhyXMiniProblems.ProblemPhyXMini0428.MatchesProblemStatement}
  \uses{def:physics:phyx-mini-0428:phyxminiproblems-problemphyxmini0428-monatomicidealgasheatenginecycle}
  Prose-level model assumptions: a closed monatomic ideal-gas heat engine. The value 5/3 is the heat-capacity ratio of a monatomic ideal gas
\end{definition}
\begin{proof}
  This is the declared model definition of Matches Problem Statement.
\end{proof}
\begin{definition}[Matches Supplied Pressure Volume Diagram]
  \label{def:physics:phyx-mini-0428:phyxminiproblems-problemphyxmini0428-matchessuppliedpressurevolumediagram}
  \lean{PhyXMiniProblems.ProblemPhyXMini0428.MatchesSuppliedPressureVolumeDiagram}
  \uses{def:physics:phyx-mini-0428:phyxminiproblems-problemphyxmini0428-pressureinpascals, def:physics:phyx-mini-0428:phyxminiproblems-problemphyxmini0428-pressureinkilopascals, def:physics:phyx-mini-0428:phyxminiproblems-problemphyxmini0428-volumeincubiccentimetres, def:physics:phyx-mini-0428:phyxminiproblems-problemphyxmini0428-temperatureinkelvins, def:physics:phyx-mini-0428:phyxminiproblems-problemphyxmini0428-statelabel, def:physics:phyx-mini-0428:phyxminiproblems-problemphyxmini0428-cycleleg, def:physics:phyx-mini-0428:phyxminiproblems-problemphyxmini0428-diagramlabel, def:physics:phyx-mini-0428:phyxminiproblems-problemphyxmini0428-expectedlegstart, def:physics:phyx-mini-0428:phyxminiproblems-problemphyxmini0428-expectedlegfinish, def:physics:phyx-mini-0428:phyxminiproblems-problemphyxmini0428-expectedprocesskind, def:physics:phyx-mini-0428:phyxminiproblems-problemphyxmini0428-expectedpathgeometry, def:physics:phyx-mini-0428:phyxminiproblems-problemphyxmini0428-monatomicidealgasheatenginecycle}
  Primary-image evidence: axes and units, visible labels, directed path types, the three volume coordinates, the 100 kPa pressure at state 3, the common pressure p\_max at states 1 and 2, and the 600 K readout at state 2. No heat, work, or efficiency conclusion occurs here
\end{definition}
\begin{proof}
  This is the declared model definition of Matches Supplied Pressure Volume Diagram.
\end{proof}
\begin{definition}[Has Physical Thermodynamic Parameters]
  \label{def:physics:phyx-mini-0428:phyxminiproblems-problemphyxmini0428-hasphysicalthermodynamicparameters}
  \lean{PhyXMiniProblems.ProblemPhyXMini0428.HasPhysicalThermodynamicParameters}
  \uses{def:physics:phyx-mini-0428:phyxminiproblems-problemphyxmini0428-pressureinpascals, def:physics:phyx-mini-0428:phyxminiproblems-problemphyxmini0428-volumeincubicmetres, def:physics:phyx-mini-0428:phyxminiproblems-problemphyxmini0428-temperatureinkelvins, def:physics:phyx-mini-0428:phyxminiproblems-problemphyxmini0428-statelabel, def:physics:phyx-mini-0428:phyxminiproblems-problemphyxmini0428-monatomicidealgasheatenginecycle}
  Positivity and nondegeneracy conditions selecting physical state readouts
\end{definition}
\begin{proof}
  This is the declared model definition of Has Physical Thermodynamic Parameters.
\end{proof}
\begin{definition}[Obeys Monatomic Ideal Gas Cycle Laws]
  \label{def:physics:phyx-mini-0428:phyxminiproblems-problemphyxmini0428-obeysmonatomicidealgascyclelaws}
  \lean{PhyXMiniProblems.ProblemPhyXMini0428.ObeysMonatomicIdealGasCycleLaws}
  \uses{def:physics:phyx-mini-0428:phyxminiproblems-problemphyxmini0428-pressureinpascals, def:physics:phyx-mini-0428:phyxminiproblems-problemphyxmini0428-volumeincubicmetres, def:physics:phyx-mini-0428:phyxminiproblems-problemphyxmini0428-energyinjoules, def:physics:phyx-mini-0428:phyxminiproblems-problemphyxmini0428-temperatureinkelvins, def:physics:phyx-mini-0428:phyxminiproblems-problemphyxmini0428-statelabel, def:physics:phyx-mini-0428:phyxminiproblems-problemphyxmini0428-cycleleg, def:physics:phyx-mini-0428:phyxminiproblems-problemphyxmini0428-monatomicidealgasheatenginecycle}
  Macroscopic thermodynamic laws used to analyze the cycle: pV = nRT at each equilibrium state; ΔU = (3/2)nRΔT for a monatomic ideal gas; Q\_into = ΔU + W\_by on every leg; the usual quasistatic boundary-work laws on isobaric and isochoric legs; p\_f/p\_i = (V\_i/V\_f)\textasciicircum{}γ and zero heat on an adiabatic leg. The laws are uniform in the leg and contain no numerical efficiency or answer choice
\end{definition}
\begin{proof}
  This is the declared model definition of Obeys Monatomic Ideal Gas Cycle Laws.
\end{proof}
\begin{definition}[net Work Done By Gas In Joules]
  \label{def:physics:phyx-mini-0428:phyxminiproblems-problemphyxmini0428-networkdonebygasinjoules}
  \lean{PhyXMiniProblems.ProblemPhyXMini0428.netWorkDoneByGasInJoules}
  \uses{def:physics:phyx-mini-0428:phyxminiproblems-problemphyxmini0428-energyinjoules, def:physics:phyx-mini-0428:phyxminiproblems-problemphyxmini0428-monatomicidealgasheatenginecycle}
  Net work done by the gas in one traversal of the closed cycle
\end{definition}
\begin{proof}
  This is the declared model definition of net Work Done By Gas In Joules.
\end{proof}
\begin{definition}[total Heat Input In Joules]
  \label{def:physics:phyx-mini-0428:phyxminiproblems-problemphyxmini0428-totalheatinputinjoules}
  \lean{PhyXMiniProblems.ProblemPhyXMini0428.totalHeatInputInJoules}
  \uses{def:physics:phyx-mini-0428:phyxminiproblems-problemphyxmini0428-energyinjoules, def:physics:phyx-mini-0428:phyxminiproblems-problemphyxmini0428-monatomicidealgasheatenginecycle}
  Total heat absorbed by the gas, the sum of positive signed leg heats
\end{definition}
\begin{proof}
  This is the declared model definition of total Heat Input In Joules.
\end{proof}
\begin{definition}[thermal Efficiency]
  \label{def:physics:phyx-mini-0428:phyxminiproblems-problemphyxmini0428-thermalefficiency}
  \lean{PhyXMiniProblems.ProblemPhyXMini0428.thermalEfficiency}
  \uses{def:physics:phyx-mini-0428:phyxminiproblems-problemphyxmini0428-monatomicidealgasheatenginecycle, def:physics:phyx-mini-0428:phyxminiproblems-problemphyxmini0428-networkdonebygasinjoules, def:physics:phyx-mini-0428:phyxminiproblems-problemphyxmini0428-totalheatinputinjoules}
  Dimensionless thermal efficiency W\_net / Q\_in
\end{definition}
\begin{proof}
  This is the declared model definition of thermal Efficiency.
\end{proof}
\begin{definition}[Answer Choice]
  \label{def:physics:phyx-mini-0428:phyxminiproblems-problemphyxmini0428-answerchoice}
  \lean{PhyXMiniProblems.ProblemPhyXMini0428.AnswerChoice}
  Labels of the four choices printed in the source problem
\end{definition}
\begin{proof}
  This is the declared model definition of Answer Choice.
\end{proof}
\begin{definition}[displayed Efficiency]
  \label{def:physics:phyx-mini-0428:phyxminiproblems-problemphyxmini0428-displayedefficiency}
  \lean{PhyXMiniProblems.ProblemPhyXMini0428.displayedEfficiency}
  \uses{def:physics:phyx-mini-0428:phyxminiproblems-problemphyxmini0428-answerchoice}
  Dimensionless decimal value printed beside each answer label
\end{definition}
\begin{proof}
  This is the declared model definition of displayed Efficiency.
\end{proof}
\begin{definition}[recorded Answer Choice]
  \label{def:physics:phyx-mini-0428:phyxminiproblems-problemphyxmini0428-recordedanswerchoice}
  \lean{PhyXMiniProblems.ProblemPhyXMini0428.recordedAnswerChoice}
  \uses{def:physics:phyx-mini-0428:phyxminiproblems-problemphyxmini0428-answerchoice}
  Answer label recorded by the supplied dataset
\end{definition}
\begin{proof}
  This is the declared model definition of recorded Answer Choice.
\end{proof}
\begin{definition}[Rounds To Two Decimal Places]
  \label{def:physics:phyx-mini-0428:phyxminiproblems-problemphyxmini0428-roundstotwodecimalplaces}
  \lean{PhyXMiniProblems.ProblemPhyXMini0428.RoundsToTwoDecimalPlaces}
  A dimensionless value rounds to a specified two-decimal display
\end{definition}
\begin{proof}
  This is the declared model definition of Rounds To Two Decimal Places.
\end{proof}
\begin{definition}[exact Efficiency From Figure]
  \label{def:physics:phyx-mini-0428:phyxminiproblems-problemphyxmini0428-exactefficiencyfromfigure}
  \lean{PhyXMiniProblems.ProblemPhyXMini0428.exactEfficiencyFromFigure}
  Exact dimensionless efficiency predicted by the figure and laws
\end{definition}
\begin{proof}
  This is the declared model definition of exact Efficiency From Figure.
\end{proof}
\begin{lemma}[derived Temperature Heat And Work Readouts]
  \label{lem:physics:phyx-mini-0428:phyxminiproblems-problemphyxmini0428-derivedtemperatureheatandworkreadouts}
  \lean{PhyXMiniProblems.ProblemPhyXMini0428.derivedTemperatureHeatAndWorkReadouts}
  \uses{def:physics:phyx-mini-0428:phyxminiproblems-problemphyxmini0428-energyinjoules, def:physics:phyx-mini-0428:phyxminiproblems-problemphyxmini0428-temperatureinkelvins, def:physics:phyx-mini-0428:phyxminiproblems-problemphyxmini0428-monatomicidealgasheatenginecycle, def:physics:phyx-mini-0428:phyxminiproblems-problemphyxmini0428-matchesproblemstatement, def:physics:phyx-mini-0428:phyxminiproblems-problemphyxmini0428-matchessuppliedpressurevolumediagram, def:physics:phyx-mini-0428:phyxminiproblems-problemphyxmini0428-hasphysicalthermodynamicparameters, def:physics:phyx-mini-0428:phyxminiproblems-problemphyxmini0428-obeysmonatomicidealgascyclelaws, def:physics:phyx-mini-0428:phyxminiproblems-problemphyxmini0428-networkdonebygasinjoules}
  The figure and governing laws first determine T₁ = 100 K and T₃ = 100 / 6\textasciicircum{}(2/3) K. The first two leg heats and the total work then follow from the monatomic internal-energy law and the first law. These are derived conclusions, not fields of a scenario, figure, or law structure
\end{lemma}
\begin{proof}
  The conclusion follows from the governing relations and figure readouts named in the statement of derived Temperature Heat And Work Readouts.
\end{proof}
% --- Archon declaration topology end ---
% --- Archon physics formalization source end ---
